% Part: second-order-logic
% Chapter: syntax-and-semantics
% Section: semantic-notions

\documentclass[../../../include/open-logic-section]{subfiles}

\begin{document}

\olfileid{sol}{syn}{sem}
\olsection{语义概念}

\begin{explain}
\emph{有效性}、\emph{语义后承}和\emph{可满足性}这些核心逻辑概念在二阶逻辑中的定义方式与一阶逻辑中相同，只是底层的满足关系现在是针对二阶公式的。当然，二阶语句是指所有变元（包括谓词变元和函数变元）都被约束的公式。
\end{explain}

\begin{defn}[有效性]
语句 $!A$ 是\emph{有效的}，$\Entails !A$，当且仅当对每个结构~$\Struct M$ 都有 $\Sat{M}{!A}$。
\end{defn}

\begin{defn}[语义后承]
语句集合~$\Gamma$ \emph{语义地后承}语句~$!A$，$\Gamma \Entails !A$，当且仅当对于每个满足 $\Sat{M}{\Gamma}$ 的结构~$\Struct M$，都有 $\Sat{M}{!A}$。
\end{defn}

\begin{defn}[可满足性]
语句集合~$\Gamma$ 是\emph{可满足的}，若存在某个结构~$\Struct M$ 使得 $\Sat{M}{\Gamma}$。若 $\Gamma$ 不可满足，则称其为\emph{不可满足的}。
\end{defn}

\end{document}